% Shared presentation style, adapted from reference-slides/.
\usetheme[progressbar=frametitle,numbering=fraction]{metropolis}

\usepackage[T1]{fontenc}
\usepackage{amsmath,amssymb}
\usepackage{booktabs}
\usepackage{braket}
\usepackage{microtype}
\usepackage{tabularx}
\usepackage{tikz}
\usepackage{qcircuit}
% Classical gate from the paper; this qcircuit version provides only \gate.
\providecommand{\cgate}[1]{*+<.6em>{#1} \POS ="i","i"+UR;"i"+UL
  **\dir{-};"i"+DL **\dir{-};"i"+DR **\dir{-};"i"+UR **\dir{-},"i" \cw}
\usepackage{pgfpages}
\usetikzlibrary{arrows.meta,calc,decorations.pathreplacing,fit,positioning}

% Notes stay in the source; `make notes` places them on a second screen.
\ifdefined\showspeakernotes
  \setbeameroption{show notes on second screen=right}
  % Avoid duplicate PDF destinations when pgfpages composes the two screens.
  \hypersetup{draft}
\else
  \setbeameroption{hide notes}
\fi

% Palette and typography match the reference presentation.
\definecolor{CVQCInk}{HTML}{17324D}
\definecolor{CVQCTeal}{HTML}{087E8B}
\definecolor{CVQCPale}{HTML}{E8F3F4}
\definecolor{CVQCGold}{HTML}{E5A23B}
\definecolor{CVQCGoldPale}{HTML}{FFF2D8}
\definecolor{CVQCCoral}{HTML}{C85151}
\definecolor{CVQCCoralPale}{HTML}{F9E8E6}
\definecolor{CVQCMuted}{HTML}{5E6D78}
\definecolor{CVQCFog}{HTML}{F4F7F8}

\setbeamercolor{normal text}{fg=CVQCInk,bg=white}
\setbeamercolor{alerted text}{fg=CVQCCoral}
\setbeamercolor{example text}{fg=CVQCTeal}
\setbeamercolor{frametitle}{fg=white,bg=CVQCInk}
\setbeamercolor{progress bar}{fg=CVQCTeal,bg=CVQCPale}
\setbeamercolor{progress bar in head/foot}{fg=CVQCTeal,bg=CVQCPale}
\setbeamercolor{title separator}{fg=CVQCTeal}
\setbeamercolor{palette primary}{fg=white,bg=CVQCInk}
\setbeamercolor{block title}{fg=white,bg=CVQCInk}
\setbeamercolor{block body}{fg=CVQCInk,bg=CVQCFog}
\setbeamercolor{block title alerted}{fg=white,bg=CVQCCoral}
\setbeamercolor{block body alerted}{fg=CVQCInk,bg=CVQCCoralPale}
\setbeamercolor{block title example}{fg=white,bg=CVQCTeal}
\setbeamercolor{block body example}{fg=CVQCInk,bg=CVQCPale}
\setbeamercolor{takeaway}{fg=CVQCInk,bg=CVQCPale}
\setbeamercolor{question}{fg=CVQCInk,bg=CVQCGoldPale}

\setbeamerfont{title}{series=\bfseries,size=\huge}
\setbeamerfont{frametitle}{series=\bfseries,size=\Large}
\setbeamerfont{block title}{series=\bfseries}
\setbeamertemplate{navigation symbols}{}
\metroset{block=fill}

% An atom marks a quantum party, as in Hung's QCrypt 2020 talk, p. 10:
% https://2020.qcrypt.net/slides/QCrypt2020TalkSlides058Hung.pdf
% Draw inside the party's scope so the marker follows its diagram scale.
\newcommand{\quantumpartymark}[1]{%
  \begin{scope}[shift={($(#1.north east)+(0.15,0.3)$)},
      draw=CVQCTeal,line width=1.1pt]
    \foreach \angle in {0,60,120} {
      \draw[rotate=\angle] (0,0) ellipse[x radius=0.58cm,y radius=0.21cm];
    }
    \fill[CVQCGold] (0,0) circle[radius=0.075cm];
    \fill[CVQCTeal] (0.58,0) circle[radius=0.045cm];
  \end{scope}%
}

% Reusable callouts for later slides.
\newcommand{\takeaway}[1]{%
  \begin{beamercolorbox}[wd=\textwidth,sep=0.85ex,rounded=true,center]{takeaway}
    #1
  \end{beamercolorbox}%
}
\newcommand{\questionbox}[1]{%
  \begin{beamercolorbox}[wd=\textwidth,sep=1ex,rounded=true,center]{question}
    #1
  \end{beamercolorbox}%
}
